% !TeX root = ../papers/nc-p3.tex
\section{二次判別式の分岐被覆による非平方型の全面排除}
\label{sec:nonsquare-opposite-discriminant}
\begin{theorem}\label{thm:nonsquare-opposite-discriminant}
四辺有限な標準有限Frobenius窓の異なる反対根で、対称二次式行列 $N$ の一般階数を3とする。
その余因子は四本の独立な三次式 $f_i$ により
$\operatorname{adj}N=(f_i f_j)$ と書ける。特に原始核が二次式であれば余因子の共通二次因子は平方である。
点の台、局所埋込次元、重複度による例外はない。
\end{theorem}

\begin{proof}[局所二次形式と切断長]
非平方の二次式 $h$ と原始な四本の二次式 $g_i$ により $\operatorname{adj}N=hgg^{\mathsf t}$ と仮定する。
$g_i$ は定理\ref{thm:ternary-kernel-cubic-exclusion}により独立である。
点配置 $X$ の各閉点 $p$ で、完備局所座標環を $R$、$D=\mathcal O_{X,p}$、長さを $n$、
$J=(g_i)D$、$j=\operatorname{length}(D/J)$ とする。既出の三次評価から
$\sum n=20,\ \sum j\le8$ である。$D$ はArtin完全交叉であり、
$\operatorname{length}(\operatorname{ann}_D J)=j$ である。

出力の定数基底を取り $N(p)$ を第一成分だけが1の対角行列にする。第一対角の単元を $a$、
残る対称三行のSchur補を $S$ とする。$S$ の成分は $D$ 上零である。
$g'$ を核の最後の三座標とすると、第一座標もそれらの $R$ 係数の和で、
$\operatorname{adj}S=(h/a)g'g'^{\mathsf t}$ となる。$g'$ の成分は共通の素因子を持たない。
元の原始性は局所化・完備化で高さ2以上のイデアルの高さが下がらないことによる。
$h$ は平方因子を持たないので $S$ の全成分も共通の素因子を持たない。

$S$ の成分が全て $p$ で三次以上になることはない。そうなれば $p=(1,0,0,0)$ と置き、
$x_0=1$ のチャートで $N$ の右下三行は、第一行の非対角成分の一次部 $L$ により
$L L^{\mathsf t}$ となる。元の成分が二次以下なのでこれは斉次化しても恒等式である。
残りの列が $L$ に比例すれば一般階数は2以下、比例しなければ核の第一座標は恒等零である。
いずれも仮定に反する。

$h(p)\ne0$ ならその平方根も $a$ の平方根も $R$ にある。
一般階数2の三変数形式 $S$ の二次元非退化部分は負の主二次小行列式が平方のため分裂する。
$R$ は一意分解環で、原始性とGaussの補題により
$S$ は二つの原始列 $A,B$ の対称積の単元倍になる。このとき $g'$ は $A\times B$ の単元倍である。
$D$ 上では $A_i B_j=-A_j B_i$ なので全ての $A_i A_j B_k$ および
$B_i B_j A_k$ が零である。実際添字を入れ替え一巡すると元の積の負になる。
従って $J^2=0$ で $2j\ge n$ を得る。

$h=0$ が滑らかな点では $t^2=h$ を添加した環は正則完備局所環である。
ここでも $S$ を分解できる。二重被覆は有限平坦で成分の共通因子を新たに生じない。
被覆対合は二因子を交換する。実際二因子の順序による外積（積の単元係数も掛ける）は
$t g'$ の基礎環の単元倍で、符号が反転する。したがって一因子の列を $A$ と取れば
他方にその共役を取れ、係数は基礎環の単元となる。
$A=b+t c$ と書くと $S$ は $bb^{\mathsf t}-hcc^{\mathsf t}$ の単元倍、
$g'$ は $b\times c$ の単元倍である。$D$ 上で $b_i b_j=h c_i c_j$ を
二つの外積成分の積に代入すると零になる。再び $2j\ge n$ となる。
特に滑らかな射影二次曲面なら全点で足すだけで矛盾する。
\end{proof}

\begin{proof}[二平面の交線]
$h$ が二平面の積の場合、その交線上で $h=xy$（単元平方は吸収）と置く。
$R$ から $x=u^2,y=v^2$ とした正則完備局所環へ移る。二つの符号反転のそれぞれは
$uv$ を反転するので二因子を交換し、同時反転は各因子直線を保つ。
原始列 $A$ は同時反転で偶または奇に選べる。一般の原始生成列で反転の単元係数の定数項を
$\epsilon=\pm1$ とし、符号 $\epsilon$ の射影を取れば元の列の単元倍となるためである。
もう一方の列を $u$ 反転による共役に取ったときの単元係数は両反転に不変である。
偶なら $A=b+uv c$ として先の平方零の証明が通る。奇なら $A=u b+v c$ であり、
$S$ は $x b b^{\mathsf t}-y c c^{\mathsf t}$ の単元倍、$J$ は $b\times c$ の成分で生成される。
$b,c$ の全成分が極大イデアルに入れば $S$ は三次以上となり既に排除した。
ある $b_i$ が単元としてよい（逆なら $x,y$ を交換）。$D$ 上
$x b_i^2=y c_i^2$ より $(x,y)D=yD$、また外積成分の積にこの対称行列等式を使うと
$y J^2=0$ となる。$J\xrightarrow{y}\operatorname{ann}J$ の核は
$\operatorname{ann}(x,y)$ に含まれ、その長さは $d=\operatorname{length}D/(x,y)$ である。
従って $n-j\le j+d$。交線と $X$ の全切断長は高々3なので、
これと交線外の評価から $20\le2\cdot8+3$ が必要となり矛盾する。
\end{proof}

\begin{proof}[二次錐の頂点での分解]
$h$ が階数3なら頂点だけを扱えばよい。その完備座標で
$h=z^2-xy$ とし、四変数冪級数環へ
\[
 x=u_0v_0,\quad y=u_1v_1,\quad
 z=(u_0v_1+u_1v_0)/2,\quad t=(u_0v_1-u_1v_0)/2
\]
で移す。$u$ に重み1、$v$ に重み $-1$ を与えた不変環は
$R\oplus tR$, $t^2=h$ で、交換 $u\leftrightarrow v$ はその対合である。
これは等しい個数の $u,v$ を持つ単項式を四つの積で表し、唯一の行列式関係を用いればよい。
$S$ の成分の原始性は移した後も保たれる。実際代数的な座標上で四つの積による
階数一行列への射は原点外で繊維次元1、原点繊維は二平面の和で次元2である。
三次元の二重被覆は元の三空間上有限であり、余次元2以上の零点集合の逆像は
四空間で次元2以下である。局所化・完備化も余次元を保つ。ここで $S$ の単元分母は払ってよい。

この正則環での二因子の一方 $A$ を単一のトーラス重み $k$ の原始列に取れる。
トーラスは連結なので各因子直線を保つ。原始列の最低非零次数部はこの直線への作用の
一つの文字を持ち、その重みへの射影は列の単元倍になる。
実際原始列の生成する部分加群は完備位相で閉じており、各有限ジェットのトーラス表現で
重み射影がこの部分加群の像を保つ。射影は最低部を変えない。
交換は $t$ の符号から因子を交換するので、もう一方を $A(v,u)$ と取れば単元係数は不変かつ交換不変で
$R$ に属する。必要なら交換して $k\ge0$ とする。重み $k$ の項の全次数は少なくとも $k$。
よって $S$ の $R$ での位数は少なくとも $k$ となり、三次接触の排除から $k\le2$ である。

$k=0$ は $A=b+t c$ で既述の平方零の場合である。$k=1$ なら $A=u_0 b+u_1 c$ と
$b,c$ を $R$ 上に取れる。不変係数の $t$ は $tu_0=zu_0-xu_1$、$tu_1=yu_0-zu_1$
で吸収する。このとき $S$ は
$x b b^{\mathsf t}+z(bc^{\mathsf t}+cb^{\mathsf t})+y c c^{\mathsf t}$ の単元倍、
$J$ は $b\times c$ の成分で生成される。全成分が極大イデアルに入れないので、
交換および局所的な行基底変更で $b=(1,0,0)^{\mathsf t}$ とする。
$c=(\beta,c_2,c_3)^{\mathsf t}$ なら $J=(c_2,c_3)D$、
$x+2\beta z+\beta^2 y=0$、$(z+\beta y)J=0$、$yJ^2=0$ である。
$J$ での $y$ 倍の核は $(x,y,z)$ で消えて長さ1以下だから $n\le2j+1$ を得る。

$k=2$ のとき $A$ の全次数2の部は $u_0^2,u_0u_1,u_1^2$ の定数列係数の和である。
三係数が独立なら $S$ の二次部は $x^2,xz,2z^2-xy,xy,yz,y^2$ を張る。
従って $D$ の極大イデアルの二乗は零であり、Gorenstein性から $n\le2$。
また $g'$ は $A(u,v)\times A(v,u)$ を $t$ で割った単元倍なので
$p$ で零であり $j\ge1$ である。従って $n\le2j$。
三係数が従属なら $S$ の二次部には定数核がある。この場合を次に排除する。
\end{proof}

\begin{proof}[二次接触に共通の定数核がある場合]
一般に階数3の $h$ の頂点で $S$ が二次以上であり、その二次部に定数核があるとする。
定数基底でその最後の行列を零にし、$p=(1,0,0,0)$、$w=(x_1,x_2,x_3)$ とする。
$N$ の右下ブロックは $H=L L^{\mathsf t}+S_2$ で、$L$ は三つの一次式、
$S_2$ は最後の行列が零の $w$ の二次式行列である。
第一成分とその下の列は $x_0^2+\alpha x_0+f$ と $x_0L+b$ と書ける。
$\det H=h g_0^2\ne0$ より $L_3\ne0$ で $H^{-1}L=e_3/L_3$ である。
$\det N=0$ の $x_0$ の係数を読むと $b_3=\alpha L_3/2$。
従って $v=(x_0+\alpha/2,L^{\mathsf t})^{\mathsf t}$ と置いて
\[
 N=vv^{\mathsf t}+\begin{pmatrix}H'&0\\0&0\end{pmatrix}
\]
を得る。$H'$ は $w$ だけの三行の対称二次式行列である。
$N$ の第四成分で消去すれば $H'$ は一般階数2、その余因子と元の余因子の先頭三成分には
$L_3^2\operatorname{adj}H'=h(g_0,g_1,g_2)^{\mathsf t}(g_0,g_1,g_2)$ の関係がある。
$h$ は既約なので各 $g_i$ は $L_3$ の倍である。商 $d_i$ はその平方が
$x_0$ に依存しないので $w$ の一次式であり、三本は元の独立性により独立である。

$d=(d_0,d_1,d_2)^{\mathsf t}$ として $H'd=0$ より
$H'=[d]_\times C[d]_\times^{\mathsf t}$ と定数対称行列 $C$ を使って書ける。
ここで $[d]_\times$ は $d$ との外積の行列。この表示を確認するには
$F(d,b)=b^{\mathsf t}H'(d)b$ を取る。これは $b\mapsto b+td$ に不変で
$(d,b)$ の二重次数が $(2,2)$。従って二行を変換する $\mathrm{SL}_2$ の最高重み0のベクトルであり、
標数零の有限次元表現の完全可約性から不変である。
階数2の二行行列は $d\times b$ が同じなら同一軌道で、どの一小行列が非零のチャートでも
その三成分から正則に代表を選べる。三空間の原点を越えて延長すれば $F$ は三小行列の二次式である。
この表示の余因子の係数は $d^{\mathsf t}(\operatorname{adj} C)d=h$ である
（対角 $C$ で直接計算し、合同変換で一般の場合を得る）。
よって $C$ は可逆で、$H'$ の六成分は $w$ の二次式を全て張る。
実際 $C=I$ に帰着すれば $(d^{\mathsf t}d)I-dd^{\mathsf t}$ の六成分である。

$N$ の第四行と第四列を使う二次小行列式より、
$X$ 上で $L_3^2$ は $p$ のイデアルの二乗に消される。従ってこの非零な二次式の
生成する巡回加群の長さは高々4である。しかし有限辺の完全跡と相手の二次評価は
そのような二次式を許さない（定理\ref{thm:nonreduced-opposite-root-separation}の最初の証明）。
これで定数核の場合も排除した。錐頂点での残りは $n\le2j+1$ を満たし、
他点での $2j\ge n$ と合わせると $20\le17$ が必要になり矛盾する。
定数 $h$ または平方の $h$ の残る両場合に、その平方根を核へ掛ければ主張を得る。
\end{proof}
